% =====================================================================
% Intelligenza artificiale in psicopatologia
% Perspective per Frontiers in Digital Health -- bozza italiana
% Marco Cremaschi, giugno 2026
% Compilare con: pdflatex + bibtex + pdflatex (x2)
% =====================================================================
\documentclass[utf8]{FrontiersinHarvard}

\usepackage[T1]{fontenc}
\usepackage[italian]{babel}
\usepackage{url,lineno,microtype}
\usepackage[hyperfootnotes=false,hidelinks]{hyperref}
\usepackage[onehalfspacing]{setspace}
\usepackage{booktabs}
\usepackage{tabularx}
\usepackage{ragged2e}
\setlength{\marginparwidth}{2cm}
\usepackage[colorinlistoftodos,textsize=footnotesize]{todonotes}
\usepackage[nohyperlinks,nolist]{acronym}

\linenumbers

\def\keyFont{\fontsize{8}{11}\helveticabold}
\def\firstAuthorLast{Cremaschi and Pierotti}
\def\Authors{Marco Cremaschi\,$^{1,*}$ and Andrea Primo Pierotti\,$^{2}$}
\def\Address{$^{1}$Dipartimento di Informatica, Sistemistica e Comunicazione (DISCo), Universit\`a degli Studi di Milano-Bicocca, Milano, Italia \\
$^{2}$ASST Grande Ospedale Metropolitano Niguarda, Milano, Italia}
\def\corrAuthor{Marco Cremaschi}
\def\corrEmail{\texttt{info@whattadata.it}\\ORCID: \href{https://orcid.org/0000-0001-7840-6228}{0000-0001-7840-6228}}

\newcolumntype{Y}{>{\RaggedRight\arraybackslash}X}

%TC:ignore
\newacro{AIAct}[AI Act]{Artificial Intelligence Act}
\newacro{AUC}[AUC]{area under the curve}
\newacro{CBT}[CBT]{cognitive behavioral therapy}
\newacro{FDA}[FDA]{Food and Drug Administration}
\newacro{IA}[IA]{intelligenza artificiale}
\newacro{LLM}[LLM]{large language model}
\newacroplural{LLM}[LLM]{large language models}
\newacro{ML}[ML]{machine learning}
\newacro{SaMD}[SaMD]{software as a medical device}
\newacro{WHO}[WHO]{World Health Organization}
%TC:endignore

\begin{document}
\onecolumn
\firstpage{1}

\title[IA e psicopatologia]{Intelligenza artificiale in psicopatologia: epistemologia clinica, relazione terapeutica e governance}

\author[\firstAuthorLast]{\Authors}
\address{}
\correspondance{}
\extraAuth{ORCID Andrea Primo Pierotti: \href{https://orcid.org/0000-0002-7124-5668}{0000-0002-7124-5668}}

\maketitle

%TC:ignore
\noindent\textbf{Conteggio parole, bibliografia esclusa: 1978}
%TC:endignore

%TC:ignore
\todo[inline]{Prima della sottomissione a Frontiers in Digital Health come Perspective: tradurre integralmente in inglese, confermare i contributi autore e verificare DOI, PMID, ISBN o URL funzionale.}
\todo[inline]{Cercare e integrare lavori simili o prospettive affini su IA, psicopatologia, diagnosi relazionale, chatbot terapeutici e governance clinica, distinguendo contributi concettuali, revisioni e studi empirici.}
%TC:endignore

\begin{abstract}
\section{}
L'\ac{IA} sta entrando nella salute mentale attraverso chatbot, strumenti di intake, modelli predittivi, analisi del linguaggio, digital phenotyping, documentazione automatica e sistemi multimodali. Questa Perspective sostiene che la sua applicazione in psicologia e psichiatria non possa seguire semplicemente il percorso osservato in altri ambiti medici. In psicopatologia l'oggetto clinico \`e soggettivo, opaco, narrativo, temporale, contestuale e relazionale; la diagnosi non \`e solo classificazione di pattern, ma comprensione del significato dell'esperienza vissuta dentro una storia e una relazione di cura. Metriche come accuratezza, \ac{AUC} o F1 restano utili per valutare modelli circoscritti, ma non bastano a stabilire utilit\`a clinica, sicurezza e responsabilit\`a. Discutiamo tre implicazioni: la diagnosi psicopatologica richiede una epistemologia clinica diversa; la relazione terapeutica non \`e riducibile a una conversazione artificialmente empatica; il chatbot diventa un nuovo oggetto relazionale che pu\`o entrare nella stanza terapeutica prima del clinico. Ne deriva una direzione per il campo: governance, evidenza e supervisione devono crescere con prossimit\`a clinica, autonomia conversazionale, vulnerabilit\`a dell'utente e capacit\`a performativa dell'output.

\tiny
\keyFont{\section{Keywords:} artificial intelligence; mental health; psychopathology; clinical epistemology; diagnosis; therapeutic alliance; chatbots; AI governance.}
\end{abstract}

% =====================================================================
\section{Introduction}
\label{sec:intro}

La domanda non \`e pi\`u se l'\ac{IA} entrer\`a nella salute mentale. \`E gi\`a presente nelle ricerche online dei pazienti, nei chatbot generalisti, nelle app di auto-aiuto, negli strumenti di trascrizione, nelle piattaforme di triage e nelle soluzioni che promettono screening, monitoraggio o supporto alla cura \citep{guo2024,hua2025}.\footnote{World Health Organization. \emph{Ethics and Governance of Artificial Intelligence for Health: Guidance on Large Multi-Modal Models}. 2024. \url{https://www.who.int/publications/i/item/9789240084759}.} La questione centrale \`e quale idea di diagnosi, terapia e responsabilit\`a clinica venga incorporata in questi sistemi.

La nostra prospettiva \`e che l'applicazione dell'\ac{IA} in psicologia e psichiatria non possa essere importata senza modifiche da settori dove il compito \`e pi\`u delimitato, il dato pi\`u standardizzato e il target pi\`u verificabile, come imaging, elettrocardiografia o monitoraggio ospedaliero \citep{abramoff2018,lang2023,ribeiro2020,yao2021,adams2022}. In psicopatologia il dato clinico \`e anche un racconto, un silenzio, un tono di voce, una sequenza temporale, una relazione, una cultura e un modo di attribuire significato all'esperienza. Una diagnosi di depressione, per esempio, non diventa pienamente clinica se si limita a rilevare umore basso, sonno alterato o riduzione della mobilit\`a: deve comprendere perdita, colpa, rischio, funzionamento relazionale, risorse e risposta del paziente alla formulazione stessa.

Questa Perspective discute avanzamenti attuali e direzioni future. L'\ac{IA} pu\`o essere utile come tecnologia di osservazione, documentazione, monitoraggio, formazione e supporto al giudizio. Diventa problematica quando trasforma la psicopatologia in classificazione decontestualizzata o simula una relazione di cura senza assumere le condizioni cliniche, corporee e responsabili dell'incontro terapeutico.

% =====================================================================
\section{Una diversa epistemologia della diagnosi}
\label{sec:epistemologia}

La diagnosi psicopatologica non coincide con la classificazione di un pattern. Jaspers aveva gi\`a posto al centro della psicopatologia la descrizione rigorosa dell'esperienza soggettiva e la distinzione tra spiegare causalmente e comprendere il significato dell'esperienza vissuta \citep{jaspers1963,stanghellini2013}. La filosofia della psichiatria contemporanea ha poi criticato l'idea che i sintomi siano oggetti atomici, indipendenti dal contesto e immediatamente misurabili \citep{parnas2013,sass2011}. Kendler ha sostenuto una visione pluralistica e multilivello dei disturbi psichiatrici, difficilmente riducibile a singoli marcatori biologici \citep{kendler2011,kendler2012,kendler2016}.

Questa differenza diventa cruciale nella valutazione dei sistemi di \ac{IA}. Accuratezza, F1 e \ac{AUC} misurano la corrispondenza tra output e label in uno specifico dataset; non misurano la validit\`a della label, la qualit\`a della formulazione, il significato soggettivo dei sintomi, l'effetto dell'output sul paziente, la gestione dell'incertezza o la responsabilit\`a del clinico. La domanda non \`e solo se il modello abbia classificato correttamente, ma che cosa abbia perso, quale azione orienti e che cosa accada alla relazione di cura.

La diagnosi, inoltre, \`e un processo relazionale. La valutazione collaborativa e il Therapeutic Assessment mostrano che l'assessment pu\`o diventare intervento quando aiuta il paziente a comprendere la propria esperienza e a trasformare la domanda iniziale \citep{finn2007,finn2012,poston2010}. La diagnosi non \`e quindi solo un'etichetta: \`e negoziazione, chiarificazione e restituzione. Include anche ci\`o che il paziente suscita nel clinico, come risonanze, oscillazioni dell'alleanza, rotture e riparazioni.

% =====================================================================
\section{Segnali eterogenei e misurazione frammentata}
\label{sec:segnali}

Una risposta ai limiti dei modelli descrittivi \`e integrare segnali eterogenei: linguaggio, voce, comportamento digitale, sonno, ritmo circadiano, geolocalizzazione, wearable, questionari, cartella clinica, video e marker fisiologici. Il digital phenotyping nasce da questa promessa: osservare la psicopatologia fuori dallo studio, in modo longitudinale ed ecologico \citep{onnela2016,torous2017}. Segnali vocali, fisiologici e comportamentali possono indicare cambiamenti del sonno nel disturbo bipolare, riduzione della mobilit\`a nella depressione, isolamento, alterazioni prosodiche o segnali di ricaduta. Tuttavia, ci\`o che potrebbe migliorare la diagnostica diventa esso stesso complesso: richiede dispositivi, consenso, interoperabilit\`a, validazione esterna, manutenzione dei modelli, protezione dei minori e capacit\`a clinica di interpretare proxy non univoci.

Le soluzioni oggi disponibili riflettono questa frammentazione. Canvas Dx usa input strutturati, video e informazioni cliniche per una valutazione regolata dell'autismo pediatrico.\footnote{U.S. Food and Drug Administration. \emph{De Novo Classification Request for Cognoa ASD Diagnosis Aid, DEN200069}. 2021. \url{https://www.accessdata.fda.gov/cdrh_docs/reviews/DEN200069.pdf}.} Altri strumenti lavorano su intake, supporto conversazionale, qualit\`a della seduta, documentazione o biomarcatori vocali \citep{fitzpatrick2017,inkster2018}.\footnote{Pagine ufficiali consultate giugno 2026: Limbic Access, \url{https://limbic.ai/access}; Lyssn, \url{https://www.lyssn.io/}; Eleos Health, \url{https://eleos.health/}; Aiberry, \url{https://www.aiberry.com/}.} Il punto non \`e confrontare prodotti, ma mostrare che il campo tenta di misurare dimensioni diverse della vita mentale senza averle ancora integrate in un modello clinicamente validato, generalizzabile e governabile.

\begin{table}[htbp]
\centering\small
\caption[Dimensioni misurate e complessit\`a clinica dei segnali in psicopatologia]{Dimensioni misurate e complessit\`a clinica dei segnali in psicopatologia. Ogni fonte informativa aumenta la ricchezza osservativa ma introduce anche vincoli interpretativi, infrastrutturali e relazionali.}
\label{tab:frammentazione}
\begin{tabularx}{\textwidth}{>{\RaggedRight}p{3.0cm} Y Y}
\toprule
\textbf{Dimensione misurata} & \textbf{Promessa clinica} & \textbf{Fattore di complessit\`a} \\
\midrule
Sintomi descrittivi & Screening, triage, follow-up. & Label eterogenee, comorbidit\`a, riduzione della formulazione a checklist. \\
\addlinespace
Linguaggio e voce & Rilevare disorganizzazione, umore, rischio. & Ironia, cultura, contesto, negazione, significato relazionale della parola. \\
\addlinespace
Comportamento digitale & Monitorare sonno, mobilit\`a, ritmo e isolamento. & Sorveglianza, falsi proxy, device burden, consenso, disuguaglianze tecnologiche. \\
\addlinespace
Segnali fisiologici e wearable & Indicatori di arousal, ritmo e cambiamento. & Infrastruttura, rumore, interoperabilit\`a, interpretazione clinica. \\
\addlinespace
Conversazione generativa & Accesso, psicoeducazione, supporto tra sedute. & Sycophancy, allucinazioni, dipendenza, falsa empatia, mancata escalation. \\
\bottomrule
\end{tabularx}
\end{table}

La direzione futura non \`e ridurre questi segnali, ma renderli clinicamente interpretabili. L'\ac{IA} pu\`o sostenere diagnosi, decision support e indicazione terapeutica se l'output resta una proposta da discutere, non un verdetto. Una depressione classificata correttamente non dice ancora se il lavoro debba concentrarsi su lutto, trauma, regolazione emotiva, farmacoterapia, terapia familiare o protezione dal rischio.

% =====================================================================
\section{Relazione terapeutica e relazione con l'IA}
\label{sec:relazione}

I \acp{LLM} producono risposte linguisticamente fluide, spesso percepite come empatiche \citep{lee2024empathic,welivita2024empathetic}. Ma l'empatia terapeutica non coincide con un tono validante. Include setting, responsabilit\`a professionale, storia condivisa, confini, riparazioni, capacit\`a di tollerare ambivalenza e disponibilit\`a a non colludere.

La relazione terapeutica \`e corporea e temporale: sedute, pause, silenzi, ritorni, continuit\`a e rotture costruiscono parte del trattamento. La letteratura sull'alleanza mostra che la qualit\`a della relazione \`e associata agli esiti della psicoterapia \citep{bordin1979,fluckiger2018,norcross2011}; la ricerca psicodinamica e intersoggettiva ha sottolineato il ruolo del sapere relazionale implicito e dei momenti di incontro nel cambiamento \citep{stern1998,sander2008,safran2000}.

Il paziente non \`e soltanto un sistema complesso da misurare; \`e un sistema complesso compreso nell'incontro con un altro sistema complesso. Il terapeuta ha corpo, storia, limiti, responsabilit\`a, stile e creativit\`a. Nella prospettiva dei sistemi viventi di Louis Sander, il cambiamento emerge da adattamento, riconoscimento e co-regolazione tra sistemi in relazione \citep{sander1975,sander2008}. Un sistema artificiale pu\`o simulare empatia, coerenza e memoria conversazionale; non partecipa per\`o affettivamente nello stesso modo, non rischia qualcosa di s\'e, non ripara un enactment con responsabilit\`a personale e non possiede una personalit\`a clinica maturata nel tempo.

% =====================================================================
\section{Il chatbot come nuovo oggetto relazionale}
\label{sec:chatbot}

La trasformazione pi\`u originale potrebbe riguardare l'\ac{IA} non come strumento, ma come oggetto relazionale. Il paziente pu\`o vivere il chatbot come interlocutore, diario, consulente, amico, autorit\`a, oggetto transizionale digitale o rifugio dall'incontro umano. Un adolescente isolato pu\`o consultarlo ogni notte; una persona ansiosa pu\`o usarlo per rassicurazioni ripetute; un paziente paranoide pu\`o chiedere conferme sulle proprie interpretazioni.

Il chatbot diventa cos\`i un potenziale ``terzo nella stanza''. Il terapeuta incontra un paziente che ha gi\`a ricevuto spiegazioni, autodiagnosi, consigli, validazioni o cornici interpretative. L'\ac{IA} pu\`o modificare aspettative, linguaggio, identit\`a diagnostica e fiducia nel clinico. Per questo l'anamnesi digitale dovrebbe includere domande esplicite:

\begin{enumerate}
\setlength{\itemsep}{2pt}
\item usi chatbot quando stai male o devi decidere?
\item hai chiesto consigli su diagnosi, farmaci, relazioni o rischio suicidario?
\item l'uso riduce l'ansia o aumenta dipendenza, ruminazione, isolamento o confusione?
\item che ruolo ha quel sistema: diario, confidente, terapeuta, autorit\`a, gioco?
\end{enumerate}

Il problema \`e accentuato dalla tendenza dei modelli conversazionali a essere accomodanti. La sycophancy, cio\`e l'allineamento alle preferenze o credenze dell'utente, \`e stata descritta nei modelli addestrati con feedback umano \citep{sharma2023sycophancy}. In clinica pu\`o diventare bias relazionale: una risposta sempre validante sembra empatica, ma non introduce l'alterit\`a necessaria a differenziare, interpretare o tollerare una rottura.

I rischi emergenti vanno letti con cautela causale, ma non ignorati. Sono stati discussi feedback loop tra chatbot e sofferenza mentale, dipendenza, destabilizzazione delle credenze e amplificazione di contenuti deliranti o quasi-deliranti \citep{hudon2025,morrin2026,dohnany2026}. Il punto non \`e demonizzare i chatbot, ma riconoscere che un interlocutore artificiale che parla di salute mentale non \`e psicologicamente neutro.

% =====================================================================
\section{Una governance proporzionata alla prossimit\`a clinica}
\label{sec:governance}

Il riferimento normativo converge su un principio: l'\ac{IA} in sanit\`a deve essere governata in modo proporzionato al rischio, con trasparenza, supervisione umana, qualit\`a dei dati, robustezza, cybersecurity, documentazione e chiarezza dell'uso previsto.\footnote{World Health Organization. \emph{Ethics and Governance of Artificial Intelligence for Health}. 2024. \url{https://www.who.int/publications/i/item/9789240084759}. European Commission. \emph{Regulatory Framework on Artificial Intelligence}. \url{https://digital-strategy.ec.europa.eu/en/policies/regulatory-framework-ai}. American Psychiatric Association. \emph{Applications of Artificial Intelligence in Mental Health Care}. \url{https://www.psychiatry.org/psychiatrists/practice/artificial-intelligence/applications}.} L'\acs{AIAct} europeo adotta una logica risk-based; la \ac{FDA} ha segnalato rischi specifici dei dispositivi generativi patient-facing in salute mentale, tra cui confabulazione, contenuti inappropriati, data drift, interpretazione errata e peggioramento sintomatologico.\footnote{U.S. Food and Drug Administration. \emph{Executive Summary for the Digital Health Advisory Committee Meeting: Generative Artificial Intelligence-Enabled Digital Mental Health Medical Devices}. 2025. \url{https://www.fda.gov/media/189391/download}.}

Per la salute mentale il rischio dipende anche dalla posizione relazionale del sistema. Un software che revisiona una nota clinica \`e diverso da un agente che parla con un paziente in crisi; un reminder psicoeducativo \`e diverso da un chatbot che propone diagnosi o scelte terapeutiche. Evidenza, supervisione e governance devono crescere con quattro variabili: prossimit\`a clinica, autonomia conversazionale, vulnerabilit\`a dell'utente e capacit\`a performativa dell'output.

\begin{table}[htbp]
\centering\small
\caption[Governance proporzionata per sistemi di IA in salute mentale]{Governance proporzionata per sistemi di IA in salute mentale. Al crescere della capacit\`a dell'output di orientare identit\`a, significati e comportamenti, devono crescere evidenza clinica, supervisione, responsabilit\`a ed escalation.}
\label{tab:governance}
\begin{tabularx}{\textwidth}{>{\RaggedRight}p{3.0cm} Y Y}
\toprule
\textbf{Scenario} & \textbf{Esempi} & \textbf{Requisito minimo} \\
\midrule
Bassa prossimit\`a & Sintesi documentale, ricerca in cartella, bozze amministrative. & Audit clinico, privacy, tracciabilit\`a, correzione umana. \\
\addlinespace
Clinician-facing ad alto impatto & Suicidariet\`a, triage, supporto diagnostico, indicazione terapeutica. & Validazione esterna, performance per sottogruppi, calibrazione locale, supervisione. \\
\addlinespace
Patient-facing supervisionato & Psicoeducazione prescritta, esercizi tra sedute, monitoraggio condiviso. & Consenso, limiti espliciti, contenuti validati, escalation, protezione dei minori. \\
\addlinespace
Patient-facing autonomo e vulnerabile & Chatbot ``terapeuta'', crisi, companion, consigli su farmaci o diagnosi. & Evidenza robusta, indicazione ristretta, sorveglianza post-market, human-in-the-loop reale, responsabilit\`a esplicita. \\
\bottomrule
\end{tabularx}
\end{table}

Questa matrice non sostituisce il diritto, ma evita un errore frequente: trattare tutti gli strumenti di \ac{IA} come equivalenti. In salute mentale il rischio aumenta quando l'output non descrive soltanto il mondo, ma orienta identit\`a, significati e comportamenti.

% =====================================================================
\section{Discussion}
\label{sec:discussion}

La direzione futura non dovrebbe essere una resistenza generica all'\ac{IA}, ma una sua collocazione epistemologicamente pi\`u precisa. Gli strumenti clinician-facing sono probabilmente i candidati pi\`u maturi per un'adozione responsabile: possono organizzare dati, segnalare traiettorie, ridurre carico documentale e sostenere il ragionamento senza sostituire la formulazione clinica. Gli strumenti patient-facing richiedono soglie di evidenza pi\`u alte, perch\'e entrano direttamente nell'ecologia relazionale del paziente.

La nostra posizione \`e che l'\ac{IA} in psicopatologia debba essere valutata su utilit\`a clinica, sicurezza, impatto relazionale e governance, non solo su benchmark. La diagnosi deve restare formulazione contestuale e intersoggettiva; i segnali multimodali devono essere interpretati dentro infrastrutture e responsabilit\`a proporzionate; il chatbot deve diventare oggetto dell'anamnesi digitale. Questa Perspective non \`e una revisione sistematica: le fonti e gli esempi sono selezionati per sostenere una tesi clinica e teorica, non per stimare l'efficacia delle singole tecnologie.

L'\ac{IA} potr\`a migliorare la salute mentale se amplia l'osservazione, supporta il clinico, rende pi\`u tempestivo il monitoraggio e aiuta i servizi a intercettare rischi. Non potr\`a invece servire la psicopatologia se riduce la mente a immagine, segnale o variabile strutturata. La sfida \`e costruire una psicologia digitale che non perda il proprio oggetto: una sofferenza opaca, narrativa, temporale, relazionale e situata.

% =====================================================================
{\small
\bibliographystyle{Frontiers-Harvard}
\bibliography{references}
}

\end{document}
